\bookchapter{第六十三章：視圖、狀態與導航}

\chapterintro{單頁之玩物既已能為。此章解二「成長」之惑：如何於眾頁之間往來跳轉（導航），及數據如何在頁間流動（狀態管理）。學畢此章，則能搭「列表 詳情 編輯」經典多頁 App 之骨架矣。}

\booksection{一、狀態管理之「數據流全景」}

SwiftUI 狀態之具，非惟 @State 一端。先觀此全景之圖，明「何時用何者」：

\begin{center}
\begin{tabularx}{\linewidth}{|X|X|X|X|}
\hline
\textbf{工具} & \textbf{司何事} & \textbf{孰用之} & \textbf{典型場景} \tabularnewline
\hline
@State & 頁面私有之暫態 & 當前視圖 & 計數器、輸入框文字 \tabularnewline
@Binding & 借他人之狀態 & 子視圖改父視圖 & 自定義組件 \tabularnewline
@StateObject / @Observable & 跨頁共用之邏輯 & 視圖模型（ViewModel） & 數據載入、業務邏輯 \tabularnewline
@EnvironmentObject / .environment & 全局共用之狀態 & 整個 App & 用戶登入之態、主題設置 \tabularnewline
\hline
\end{tabularx}
\end{center}

\begin{center}\small\color{gray}SwiftUI 狀態之具全景（2026 年，@Observable 已為主流）\end{center}

\begin{lstlisting}[language=Swift]
// iOS 17+ 之新法：@Observable 宏（新項目宜用）
import Observation

@Observable
class CounterModel {                 // 視圖模型
    var count = 0
    func increment() { count += 1 }
}

struct CounterView: View {
    @State private var model = CounterModel()   // 持有模型

    var body: some View {
        VStack {
            Text("汝點了 \(model.count) 次")
            Button("點我") { model.increment() }
        }
    }
}
\end{lstlisting}

予新手直截之勸：頁面私有者用 @State，跨頁或有邏輯者用 @Observable 類。舊籍之 @StateObject/@ObservableObject 乃前代寫法，能讀懂即可，新碼悉遵上式。

\booksection{二、NavigationStack：頁面跳轉之「棧」}

iOS 16 以降，多頁導航用 NavigationStack——若一疊盤，推一入一，返則逐一彈出：

\begin{lstlisting}[language=Swift]
struct HomeView: View {
    var body: some View {
        NavigationStack {
            List {
                NavigationLink("張三之賬本") {
                    DetailView(name: "張三之賬本")   // 目標頁面
                }
                NavigationLink("李四之賬本") {
                    DetailView(name: "李四之賬本")
                }
            }
            .navigationTitle("賬本列表")   // 頂部大標題
        }
    }
}

struct DetailView: View {
    let name: String
    var body: some View {
        Text("\(name) 之詳情頁")
            .navigationTitle(name)
            .navigationBarTitleDisplayMode(.inline)
    }
}
\end{lstlisting}

NavigationLink 即「可點之跳轉連結」；navigationTitle 自生導航欄之題與返鈕。欲程式化跳轉（如點「保存」而自返）者，用 @Environment(\textbackslash{}.dismiss) 以闔當前之頁——二法相濟，九成導航之需備矣。

\booksection{三、彈窗與底部標籤：sheet 與 TabView}

除「推頁」外，尚有二種經典之交互：

\begin{lstlisting}[language=Swift]
// 1. 底部彈出半屏/全屏（錄入表單最愛）
struct HomeView: View {
    @State private var showAdd = false
    var body: some View {
        Button("記一筆") { showAdd = true }
            .sheet(isPresented: $showAdd) {
                AddRecordView()     // 彈出之錄入頁
            }
    }
}

// 2. 底部標籤欄（App 之"骨架"）
struct MainTabView: View {
    var body: some View {
        TabView {
            RecordListView()
                .tabItem { Label("賬本", systemImage: "book") }
            StatsView()
                .tabItem { Label("統計", systemImage: "chart.pie") }
            SettingsView()
                .tabItem { Label("設置", systemImage: "gear") }
        }
    }
}
\end{lstlisting}

\begin{tcbinfo}{小貼士：何時用何種？}
經驗之法：主流程以 NavigationStack 推頁（列表詳情），暫時之操作以 sheet 彈窗（新增、編輯、確認），並列數大塊以 TabView（底部標籤）。三者相合，即絕大多數 App 之骨架。

\end{tcbinfo}

\booksection{四、List：iOS 之「國民列表」}

iOS 最常用而最具「蘋果味」之組件即 List——天生帶分割線、滑動而刪、點而高亮：

\begin{lstlisting}[language=Swift]
struct RecordListView: View {
    @State private var records = ["午餐 25.5", "地鐵 6.0", "薪俸 +8000"]

    var body: some View {
        List {
            ForEach(records, id: \.self) { r in
                Text(r)
            }
            .onDelete { indexSet in          // 左滑而刪
                records.remove(atOffsets: indexSet)
            }
        }
    }
}
\end{lstlisting}

列表 + NavigationLink + sheet + TabView——合此四者，汝已能搭「賬本列表 點開詳情 點+彈錄入」完整 App 之骨架矣。次章記賬本之項，即如是搭成。

\begin{tcbpractice}{來動手：搭一「三頁骨架」}
\booksection{五、程序生命之期：前台與後台}

與安卓 Activity 生命週期相類，iOS App 亦有「生、前台、後台、亡」狀態之變。SwiftUI 中以 scenePhase 環境值察之：

\begin{lstlisting}[language=Swift]
@main
struct MyApp: App {
    @Environment(\.scenePhase) private var phase

    var body: some Scene {
        WindowGroup {
            ContentView()
        }
        .onChange(of: phase) { _, newPhase in
            switch newPhase {
            case .active:    print("返前臺，復刷新")
            case .inactive:  print("過渡中（來電/切後臺瞬間）")
            case .background: print("入後臺，存草稿")
            @unknown default: break
            }
        }
    }
}
\end{lstlisting}

典型之用：返前台則刷數據，入後台則存草稿。iOS 於後台之運行管束甚嚴——常 App 入後台不久即被懸置（凍結），勿冀於後台潛作重務；誠欲定時之務，用第 68 章本地通知，而非「後台常駐」。

作一假數據版記賬本之骨架：TabView 三標籤（賬本/統計/設置） 賬本頁以 List 陳假數據、點行跳詳情（NavigationStack） 右上「+」鈕彈錄入 sheet。不貯數據、不連網，純練骨架——骨架既成，第 65 章填以真數，如水到而渠成。

\end{tcbpractice}

\begin{tcbsummary}{本章小結}
\begin{itemize}
  \item 狀態四具：@State（頁私有）、@Binding（借用）、@Observable（邏輯）、.environment（全局）
  \item 新項以 @Observable 宏作視圖模型；舊籍之 @ObservableObject 能讀即可
  \item NavigationStack 推頁；@Environment(\textbackslash{}.dismiss) 程式化而返
  \item 主流程推頁、暫事彈 sheet、並列大塊用 TabView
  \item List 者「國民列表」也：分割線、滑動而刪，啟箱即用
\end{itemize}

\end{tcbsummary}
